\chapter{Geltungsbereich und Grundbegriffe}
\label{b13guide:ch:foundations}

Dieser Leitfaden ergänzt Band~13 um eine ausgewählte Übersicht
zentraler Definitionen, Charakterisierungen und Folgerungen über
\emph{Schranken, Infima und Suprema}. Er enthält keine Beweise
oder Beweisskizzen. Die Zusammenstellung erhebt keinen Anspruch
auf eine vollständige Liste aller bekannten oder ableitbaren Sätze.
Die Nummern mit Präfix $13H$ gehören ausschließlich zu dieser Ergänzung.

\section{Voraussetzungen und Bandgrenzen}

Durchgehend ist $(P,\leq)$ eine beliebige halbgeordnete Menge
im Sinne von Band~12; $P$ darf leer sein. Die Halbordnungsaxiome
werden als Grundlage benutzt. Allgemeine Sätze über Halbordnungen
werden hier nicht erneut entwickelt. Mengentheoretische Grundlage
ist ZF ohne zusätzliches Auswahlaxiom.

Der Gegenstand sind Teilmengen von $P$ und ihre Schranken in $P$.
Existenzannahmen beziehen sich ausschließlich auf die jeweils
genannten Infima oder Suprema. Es wird weder die Existenz aller
Paarinfima oder Paarsuprema noch die aller Infima oder Suprema
vorausgesetzt. Die Theorie der Paaroperationen gehört zu Band~14,
totale Ordnungen und Wohlordnungen zu den Bänden~15--16,
Halbverbände und Verbände zu Band~45. Vollständigkeitsaxiome,
Fixpunktsätze sowie algebraische, topologische und analytische
Zusatzstrukturen sind nicht Gegenstand dieses Leitfadens.

Allgemeine Funktionensätze sowie Injektivität, Surjektivität
und Bijektivität bleiben den Bänden~5--8 vorbehalten,
Äquivalenzrelationen Band~11. Indizierte Mengenfamilien dienen
hier lediglich zur Schreibweise von Schrankenformeln.

\section{Schranken und Beschränktheit}

\begin{definition}[Schrankenmengen]\label{b13guide:def:bounds}
Für $S\subseteq P$ heißen
\[
\begin{aligned}
 \Lo_P(S)&=\{p\in P\mid\forall x\in S\;(p\leq x)\},\\
 \Hi_P(S)&=\{p\in P\mid\forall x\in S\;(x\leq p)\}
\end{aligned}
\]
die \emph{untere} beziehungsweise \emph{obere Schrankenmenge}.
Ihre Elemente heißen untere beziehungsweise obere Schranken von $S$
in $P$. Diese Mengen entsprechen $\LB_{P,\leq}(S)$ und
$\UB_{P,\leq}(S)$ im Hauptband. Bei festem Träger und fester
Ordnung schreiben wir kurz $\Lo(S)$ und $\Hi(S)$.
\end{definition}

\begin{definition}[Beschränktheit]\label{b13guide:def:bounded}
Eine Teilmenge $S\subseteq P$ heißt \emph{nach unten beschränkt},
wenn $\Lo(S)\ne\varnothing$, und \emph{nach oben beschränkt},
wenn $\Hi(S)\ne\varnothing$. Sie heißt \emph{beschränkt}, wenn
beides gilt. \emph{Nach oben unbeschränkt} bedeutet
$\Hi(S)=\varnothing$; nach unten entsprechend $\Lo(S)=\varnothing$.
\end{definition}

Als Kurznotation verwenden wir für $p\in P$
\[
 \mathord\downarrow p=\{x\in P\mid x\leq p\},\qquad
 \mathord\uparrow p=\{x\in P\mid p\leq x\}.
\]
Die Zeichen $\min$ und $\max$ bezeichnen, wie in Band~12,
das kleinste beziehungsweise größte Element einer Menge.

\section{Infimum und Supremum}

\begin{definition}[Infimum und Supremum]\label{b13guide:def:extrema}
Ein $m\in P$ heißt \emph{Infimum} von $S\subseteq P$, wenn
\[
 m\in\Lo(S)\quad\text{und}\quad
 \forall \ell\in\Lo(S)\;(\ell\leq m).
\]
Ein $s\in P$ heißt \emph{Supremum} von $S$, wenn
\[
 s\in\Hi(S)\quad\text{und}\quad
 \forall u\in\Hi(S)\;(s\leq u).
\]
Sofern sie existieren, schreiben wir
\[
 \inf_P S=\max\Lo_P(S),\qquad
 \sup_P S=\min\Hi_P(S),
\]
oder kurz $\inf S$ und $\sup S$.
Eine Aussage wie $s=\sup S$ schließt die Existenz des Supremums ein;
bei Ungleichungen mit solchen Termen wird die Existenz ausdrücklich
vorausgesetzt. Es werden keine Ersatzwerte für nicht existierende
Infima oder Suprema eingeführt.
\end{definition}

\begin{theorem}[Eindeutigkeit und notwendige Beschränktheit]
\label{b13guide:thm:unique}
Jede Teilmenge $S\subseteq P$ besitzt höchstens ein Infimum und
höchstens ein Supremum. Ein existierendes Infimum ist eine untere
Schranke in $P$, ein existierendes Supremum eine obere Schranke in $P$.
Insbesondere gilt
\[
\begin{aligned}
 \inf S\text{ existiert}&\ \Rightarrow\ \Lo(S)\ne\varnothing,\\
 \sup S\text{ existiert}&\ \Rightarrow\ \Hi(S)\ne\varnothing.
\end{aligned}
\]
Weder das Infimum noch das Supremum muß zu $S$ gehören.
Beschränktheit allein garantiert ihre Existenz nicht.
\end{theorem}

\section{Vertauschen von oben und unten}

\begin{theorem}[Dualität der Schrankenbegriffe]
\label{b13guide:thm:dual}
Bezeichne $P^{\mathrm{op}}$ denselben Träger mit umgekehrter Ordnung,
also $x\leq_{P^{\mathrm{op}}}y$ genau dann, wenn $y\leq_P x$.
Für jedes $S\subseteq P$ gilt
\[
 \Lo_{P^{\mathrm{op}}}(S)=\Hi_P(S),\qquad
 \Hi_{P^{\mathrm{op}}}(S)=\Lo_P(S).
\]
Nach unten beschränkt in $P^{\mathrm{op}}$ ist dasselbe wie
nach oben beschränkt in $P$, und umgekehrt.
Einschließlich der jeweiligen Existenzbedingungen gilt
\[
 \inf_{P^{\mathrm{op}}}S=\sup_P S,\qquad
 \sup_{P^{\mathrm{op}}}S=\inf_P S.
\]
\end{theorem}
